% Course: 77972 - Physical Models of Biological Systems
% Academic year / semester: 2026-2027, semester A
% Actual lecture number: Unknown
% Summary/chapter index: 05; this is not an independently verified lecture number.
% Lecture date: 2026-11-17
% Summary title: Continuous Distributions, Fluorescence Compensation, and PCA
% Lecturer: Unknown
% Summarizer: Codex
% Source transcript (immutable): Transcriptions/77972 2026-11-17 שעור א [c02a1fcf-1c2d-4b79-9f2b-b4fc00680f5a]_panopto.md
% Recording: https://huji.cloud.panopto.eu/Panopto/Pages/Viewer.aspx?id=c02a1fcf-1c2d-4b79-9f2b-b4fc00680f5a
% Session ID: c02a1fcf-1c2d-4b79-9f2b-b4fc00680f5a
% Supporting source: models-notes-2026.pdf, sections 4.6-4.7 and 5.1-5.6.
% Source mapping: 00:01:48-00:22:20 ANI, marker genes, UniFrac;
% 00:22:20-00:58:45 densities, empirical CDF, entropy and scale;
% 01:00:00-01:18:23 Gaussian entropy, von Mises, biological rhythms;
% 01:18:23-01:44:44 cytometry, gating, spectra, antibody specificity;
% 02:00:02-02:24:56 spillover, compensation, arcsinh, MI calibration;
% 02:24:56-02:43:23 multivariate Gaussian, mutual information, PCA;
% 02:43:23-02:49:56 nonlinear structure and distance-based ordination.
% Editorial treatment: fill algebraic gaps; flag substantive oral ambiguities.
% Worked numerical examples and schematic figures are reconstructed, not data.
% Scope ends with PCA and the brief PCoA/MDS discussion; eigenworms postponed.

\chapter{Continuous Distributions, Fluorescence Compensation, and PCA}
\label{ch:l05-continuous}
\noindent\textit{Recording date: 17 November 2026. Summary index 05; actual lecture number unknown.}

\section*{Overview}
Genetic similarity allows comparison of microbial communities, and continuous probability distributions describe quantitative single-cell traits. In multi-channel fluorescence measurements, optical spillover mixes biological signals and creates instrument-induced covariance. Distinguishing measurement artifacts from genuine biological correlation requires calibrated matrix inversion. Covariance matrices and multivariate Gaussian models then establish the physical and mathematical foundation for principal component analysis (PCA), whitening, and dimensionality reduction.

\section*{Learning objectives}
\begin{itemize}
\item Distinguish genetic identity, phylogenetic distance, and community dissimilarity.
\item Estimate a density or cumulative distribution without confusing density with probability.
\item Explain coordinate dependence of differential entropy and invariance of mutual information under separate invertible transformations.
\item Derive Gaussian entropy and the concentrated limit of a circular distribution.
\item Derive instrument-induced covariance and invert a calibrated spillover model.
\item Separate compensation, PCA rotation, whitening, and dimension reduction, and state their physical limitations.
\end{itemize}

\section{From genome similarity to community distance}
\label{sec:l05-phylogeny}
\subsection{What does it mean for organisms to be similar?}
\label{subsec:l05-ani}
Diversity measures within a community can incorporate species similarities through a similarity matrix. While classical taxonomic ranks assign categories, genome sequence data provide quantitative metrics of divergence.

\begin{definition}{Average nucleotide identity}{l05-ani}
Average nucleotide identity (ANI) measures mean nucleotide identity across aligned homologous genomic regions of two organisms. Reported ANI values must include the fraction of the genome successfully aligned (alignment coverage).
\end{definition}
Two genomes can share close identity across a small conserved segment yet diverge across the remainder of their genomes. A sequence identity of 99\% is incomplete without alignment coverage. Operational thresholds near 95--96\% define microbial species in genomic databases, though they represent practical conventions rather than rigid natural laws. An observed gap in an identity histogram depends on the sampled taxa and alignment criteria.

For distant taxa, whole-genome alignment breaks down as alignable coverage shrinks. In such regimes, conserved marker genes provide reliable reference alignments: common panels utilize approximately 120 single-copy core genes or the 16S ribosomal RNA gene. Sequence differences across these conserved markers guide the construction of a \emph{phylogenetic tree}. Branch lengths reflect estimated sequence substitutions (evolutionary distance) rather than elapsed calendar time, unless an explicit molecular clock is established.

\begin{remark}
Marker gene panels provide practical approximations. Gene loss, duplications, horizontal gene transfer, and lineage-specific rates occur regularly. Neither universal single-copy presence nor immediate lethality following duplication holds universally. Inferred tree topologies and branch lengths carry statistical uncertainty.
\end{remark}

If $d_{ij}$ is the path length between taxa $i$ and $j$ on an inferred tree, a valid similarity matrix is
\begin{equation}
Z_{ij}=\exp(-d_{ij}/\ell),\qquad \ell>0,
\label{eq:l05-tree-similarity}
\end{equation}
where $\ell$ has the physical units of $d_{ij}$. For dimensionless distances, setting $\ell=1$ recovers $\exp(-d_{ij})$. Closely related taxa have similarities near one, while distant taxa approach zero. This similarity matrix enters taxonomic diversity indices. It represents relatedness between individual taxa, distinct from the aggregate distance between two entire communities.

\subsection{UniFrac compares represented evolutionary history}
\label{subsec:l05-unifrac}
Suppose microbial samples $A$ and $B$ contain distinct mixtures of organisms. Counting shared species treats all species pairs as equally divergent. UniFrac instead measures phylogenetic branch lengths unique to each community.

For every branch $b$ with length $L_b\geq0$, let $a_b=1$ if at least one descendant taxon occurs in sample $A$, and zero otherwise. Define $b_b$ similarly for sample $B$. Restrict evaluation to the sub-tree formed by the union of both communities.
\begin{definition}{Unweighted UniFrac}{l05-unweighted}
For a non-empty represented tree with positive total length,
\begin{equation}
d_{\mathrm{UF}}(A,B)
=\frac{\sum_b L_b\abs{a_b-b_b}}
       {\sum_b L_b\max(a_b,b_b)}.
\label{eq:l05-unifrac}
\end{equation}
This metric represents the fraction of total branch length unique to either community.
\end{definition}
\textit{Significance:} replacing a taxon with a close relative alters only a short terminal branch. Replacing it with a distant lineage alters a deep ancestral path. The metric is dimensionless and bounded in $[0, 1]$. Identical lineage sets yield zero. Completely disjoint lineages yield one. Calculations must treat the root and common ancestral branches consistently.

\begin{example}{Three-branch distance calculation}{l05-unifrac-example}
Consider three branches: length $0.3$ unique to $A$, length $0.2$ unique to $B$, and length $0.5$ shared between both communities. Then
\begin{equation}
d_{\mathrm{UF}}=\frac{0.3+0.2}{0.3+0.2+0.5}=0.5.
\label{eq:l05-unifrac-half}
\end{equation}
Half of the represented evolutionary history is unique to one community. This metric evaluates presence-absence rather than relative cell counts.
\end{example}

\begin{figure}[htbp]
\centering
\begin{tikzpicture}[x=1cm,y=1cm]
\draw[very thick,gray] (0,0)--(3,0);
\draw[very thick,primary] (3,0)--(5,1.2);
\draw[very thick,accent] (3,0)--(5,-1.2);
\fill (0,0) circle (2pt);
\node[below] at (0,0) {root};
\node[above] at (1.5,0) {$0.5$: shared};
\node[above,sloped] at (4,0.6) {$0.3$: $A$ only};
\node[below,sloped] at (4,-0.6) {$0.2$: $B$ only};
\node[right] at (5,1.2) {taxon in $A$};
\node[right] at (5,-1.2) {taxon in $B$};
\end{tikzpicture}
\caption{Schematic tree for the UniFrac calculation. Both communities share the ancestral branch ($L=0.5$), while terminal branches are unique to $A$ ($L=0.3$) and $B$ ($L=0.2$). Numerical values define the metric.}
\label{fig:l05-unifrac-tree}
\end{figure}

Unweighted UniFrac ignores cell abundance: detecting a single rare cell impacts the score as much as detecting an abundant species. Weighted UniFrac accounts for relative abundance. Let $A_b$ and $B_b$ denote the fraction of each community descending from branch $b$. A standard normalized weighted convention is
\begin{equation}
d_{\mathrm{wUF}}(A,B)=
\frac{\sum_b L_b\abs{A_b-B_b}}
     {\sum_b L_b(A_b+B_b)}.
\label{eq:l05-weighted-unifrac}
\end{equation}
The numerator weights branch lengths by abundance differences. By the triangle inequality, $\abs{A_b-B_b}\leq A_b+B_b$, bounding $d_{\mathrm{wUF}}$ in $[0, 1]$ when the denominator is positive. Evaluating weighted UniFrac requires lineage abundances across the tree.

\section{Continuous measurements: densities and cumulative probabilities}
\label{sec:l05-pdf}
\subsection{A density is probability per unit coordinate}
\label{subsec:l05-density}
Measured single-cell parameters (such as fluorescence intensity or forward scatter) behave as continuous random variables despite finite instrument resolution. Let $X$ denote a scalar random variable, $x$ an observed coordinate, $N$ the total sample size, and $\Delta N$ the observation count within a bin of width $\Delta x$ centered at $x$.
\begin{definition}{Probability density function}{l05-pdf}
For an absolutely continuous random variable, the probability density function (PDF) $p_X$ satisfies
\begin{equation}
\Pr(a\leq X\leq b)=\int_a^b p_X(x)\,\dd x,
\quad p_X(x)\geq0,\quad \int_{-\infty}^{\infty}p_X(x)\,\dd x=1.
\label{eq:l05-density-definition}
\end{equation}
For smooth densities, the empirical limit is
\begin{equation}
p_X(x)=\lim_{\Delta x\to0}\left[
\lim_{N\to\infty}\frac{\Delta N}{N\Delta x}\right],
\qquad \Pr\left(\abs{X-x}\leq\frac{\Delta x}{2}\right)
\simeq p_X(x)\Delta x.
\label{eq:l05-density-limit}
\end{equation}
\end{definition}
\textit{Significance:} probability corresponds to integrated area rather than vertical curve height. If $x$ has units of micrometers, $p_X$ has physical dimensions of $\si{\micro\metre}^{-1}$. The product $p_X\Delta x$ is dimensionless. Densities can exceed unity in narrow distributions. For continuous variables, $\Pr(X=x)=0$ at any exact point value.

The order of limits is critical. Shrinking bin width $\Delta x \to 0$ at fixed $N$ yields empty bins rather than an accurate density. Estimating a continuous density requires an asymptotic regime where $N \to \infty$ and $\Delta x \to 0$ such that the expected count $N\Delta x \to \infty$.

\subsection{Histograms, empirical distributions, and binning}
\label{subsec:l05-cdf}
\begin{definition}{Cumulative distribution function}{l05-cdf}
The cumulative distribution function (CDF) is
\begin{equation}
F_X(x)=\Pr(X\leq x)=\int_{-\infty}^x p_X(u)\,\dd u.
\label{eq:l05-cdf}
\end{equation}
The CDF is dimensionless, non-decreasing, and bounded in $[0, 1]$. Where differentiable, $F_X'(x)=p_X(x)$.
\end{definition}
For $N$ observations $x_1,\ldots,x_N$, constructing an empirical distribution requires no bin selection:
\begin{equation}
\widehat F_N(x)=\frac1N\sum_{n=1}^N\mathbf{1}_{\{x_n\leq x\}},
\qquad
\widehat p_N(x)=\frac1N\sum_{n=1}^N\delta(x-x_n).
\label{eq:l05-empirical}
\end{equation}
Here $\mathbf{1}$ denotes the indicator function and $\delta$ denotes the Dirac delta distribution. The empirical CDF steps upward by $1/N$ at each data point. Its distributional derivative equals the sum of point masses. Recovering a smooth PDF from empirical data requires smoothing or parametric model estimation.

For a histogram bin $j$ with count $n_j$ and width $\Delta x_j$, the normalized density estimate is
\begin{equation}
\widehat p_j=\frac{n_j}{N\Delta x_j},\qquad
\sum_j\widehat p_j\Delta x_j=1.
\label{eq:l05-histogram}
\end{equation}
Wide bins smooth statistical noise but obscure narrow physical peaks. Narrow bins resolve fine structure but amplify Poisson sampling noise. Validating an empirical conclusion requires testing stability across varying bin widths. Non-uniform or logarithmic binning requires dividing counts by local bin widths $\Delta x_j$.

\begin{example}{Unequal bin widths alter density heights}{l05-bins}
In a sample of $N=1000$ cells, suppose 100 cells fall in $[1, 2]$ and 100 cells fall in $[2, 4]$. Both intervals have equal empirical probability $0.1$. However, their estimated densities differ:
\begin{equation}
\widehat{p}_1 = \frac{100}{1000 \times 1} = 0.1, \qquad \widehat{p}_2 = \frac{100}{1000 \times 2} = 0.05.
\end{equation}
Plotting raw counts rather than counts divided by bin width distorts the probability density.
\end{example}

\subsubsection*{Method: comparing empirical distributions with physical models}
\begin{enumerate}
\item Define physical variables, measurement units, support bounds, and instrument resolution.
\item When plotting histograms, normalize counts by both total sample size $N$ and bin width $\Delta x$.
\item Compare candidate models directly against the empirical CDF $\widehat F_N(x)$ to avoid binning artifacts.
\item Fit model parameters using maximum likelihood estimation rather than binned least-squares.
\end{enumerate}
Cumulative distribution comparisons avoid arbitrary bin choices. However, adjacent points on an empirical CDF possess correlated errors because they accumulate overlapping sample subsets. Parameter estimation should rely on unbinned maximum likelihood estimators.

\section{Differential entropy and the scale of a coordinate}
\label{sec:l05-entropy}
\subsection{The uniform distribution exposes the issue}
\label{subsec:l05-uniform}
Let $X$ be uniform on $[x_{\min},x_{\max}]$, with width $L=x_{\max}-x_{\min}>0$:
\begin{equation}
p_X(x)=\begin{cases}1/L,&x_{\min}\leq x\leq x_{\max},\\0,&\text{otherwise}.\end{cases}
\label{eq:l05-uniform}
\end{equation}
The area is $(1/L)L=1$. A uniform density maximizes differential entropy on a fixed finite interval under no additional constraints.

\begin{definition}{Differential entropy with a reference scale}{l05-entropy}
Using natural logarithms, choose a fixed reference length $x_{\mathrm{ref}}$ in the units of $X$ and define
\begin{equation}
h_{x_{\mathrm{ref}}}(X)=-\int p_X(x)
\ln[p_X(x)x_{\mathrm{ref}}]\,\dd x.
\label{eq:l05-reference-entropy}
\end{equation}
The logarithm has a dimensionless argument. Writing $h(X)=-\int p_X\ln p_X\,\dd x$ implicitly fixes a coordinate unit or reference measure. Entropies written below without a reference subscript use this numerical-coordinate convention and are in nats.
\end{definition}
For the uniform density, direct integration gives
\begin{align}
h_{x_{\mathrm{ref}}}(X)
&=-\int_{x_{\min}}^{x_{\max}}\frac1L
      \ln\left(\frac{x_{\mathrm{ref}}}{L}\right)\dd x \notag\\
&=-\frac{L}{L}\ln\left(\frac{x_{\mathrm{ref}}}{L}\right)
=\ln\left(\frac{L}{x_{\mathrm{ref}}}\right).
\label{eq:l05-uniform-entropy}
\end{align}
At a fixed reference scale, a narrower interval has smaller differential entropy, which may be negative. Unlike discrete Shannon entropy, differential entropy is not constrained to be non-negative.

\begin{example}{Halving the interval}{l05-half-entropy}
In fixed dimensionless numerical coordinates, take $X$ uniform on $[0,1]$ and $Y=X/2$. The new interval is $[0,1/2]$, so $p_Y=2$ there. Therefore
\begin{equation}
h(X)=0,\qquad
h(Y)=-\int_0^{1/2}2\ln2\,\dd y=-\ln2.
\label{eq:l05-half-entropy}
\end{equation}
Both distributions integrate to one. The density scaling inside the logarithm accounts for the negative entropy shift.
\end{example}
\begin{remark}
A scaled interval $[0, 1/2]$ with uniform density $p_Y=2$ corresponds to $Y=X/2$. Coordinate dependence of differential entropy does not violate thermodynamic laws. A physical entropy requires an explicit reference measurement scale.
\end{remark}

\subsection{Scaling and why mutual information survives it}
\label{subsec:l05-mi-invariance}
Conservation of probability under $Y=aX$ ($a\neq0$) gives $p_Y(y)\dd y=p_X(x)\dd x$ with absolute interval widths. Hence
\begin{equation}
p_Y(y)=\frac1{\abs a}p_X(y/a).
\label{eq:l05-scale-pdf}
\end{equation}
Substitution into the entropy integral gives, in the fixed numerical-coordinate convention,
\begin{align}
h(Y)&=-\int p_X(x)\ln\left[\frac{p_X(x)}{\abs a}\right]\dd x \notag\\
&=-\int p_X(x)\ln p_X(x)\dd x
 +\ln\abs a\int p_X(x)\dd x \notag\\
&=h(X)+\ln\abs a.
\label{eq:l05-scale-entropy}
\end{align}
The same result holds with a fixed reference scale for $X$ and $Y$ in identical units. If a unit conversion scales the reference length by the same factor $a$, the normalized entropy $h_{x_{\mathrm{ref}}}$ remains invariant. Physical statements must specify which reference scale remains fixed.

\begin{proposition}{Mutual information is invariant under separate invertible reparametrizations}{l05-mi-invariance}
For continuous $X,Y$ with a joint density,
\begin{equation}
I(X;Y)=\int\!\int p_{XY}(x,y)
\ln\frac{p_{XY}(x,y)}{p_X(x)p_Y(y)}\,\dd x\,\dd y.
\label{eq:l05-continuous-mi}
\end{equation}
This quantity is non-negative, vanishes if and only if $X$ and $Y$ are independent, and remains unchanged under separate smooth invertible transformations $U=g(X)$ and $V=f(Y)$.
\end{proposition}
\begin{proof}
Let $g'$ and $f'$ be non-zero on the domains in question. The transformed joint density is $p_{UV}=p_{XY}/(\abs{g'}\abs{f'})$, while $p_U=p_X/\abs{g'}$ and $p_V=p_Y/\abs{f'}$. The Jacobians cancel in the density ratio inside the logarithm:
\begin{equation}
\frac{p_{UV}(u,v)}{p_U(u)p_V(v)} = \frac{p_{XY}(x,y)}{p_X(x)p_Y(y)}.
\end{equation}
The coordinate volume element transforms as $\dd u\,\dd v = |g'(x)f'(y)|\,\dd x\,\dd y$, canceling the Jacobian prefactor of $p_{UV}$. The integral~\eqref{eq:l05-continuous-mi} is preserved. Non-negativity follows from the non-negativity of relative entropy (Kullback-Leibler divergence).
\end{proof}
\textit{Significance:} changing the independent units or non-linear calibration of two sensors does not alter their statistical dependence. The density ratio is dimensionless.

For comparison, differential entropy transforms under coordinate change as $h(g(X))=h(X)+\mathbb E[\ln\abs{g'(X)}]$. Joint transformations that mix $X$ and $Y$ do not factor into separate coordinate reparametrizations. Such transformations alter mutual information, a property central to PCA.

\section{Gaussian widths and circular variables}
\label{sec:l05-gaussian-circle}
\subsection{The one-dimensional Gaussian and its entropy}
\label{subsec:l05-gaussian}
Fluctuating physical measurements often approach Gaussian limits when multiple independent perturbations combine. In heterogeneous cell populations, this assumption must be empirically tested. Let $\mu$ denote the mean and $\sigma>0$ the standard deviation, both in the physical units of $X$.
\begin{definition}{Gaussian density}{l05-normal}
\begin{equation}
p_X(x)=\frac1{\sqrt{2\pi}\sigma}
\exp\left[-\frac{(x-\mu)^2}{2\sigma^2}\right],
\quad \mathbb E[X]=\mu,\quad \operatorname{Var}(X)=\sigma^2.
\label{eq:l05-normal}
\end{equation}
The support covers the real line. The exponent is dimensionless, and the prefactor carries units inverse to $X$.
\end{definition}
With $z=(x-\mu)/\sigma$, normalization follows from $\int\exp(-z^2/2)\dd z=\sqrt{2\pi}$.

\begin{proposition}{Entropy measures a Gaussian's effective width}{l05-normal-entropy}
\begin{equation}
h_{x_{\mathrm{ref}}}(X)=\frac12
\ln\left(\frac{2\pi\ee\sigma^2}{x_{\mathrm{ref}}^2}\right),
\qquad
x_{\mathrm{ref}}\ee^{h_{x_{\mathrm{ref}}}(X)}=\sqrt{2\pi\ee}\,\sigma.
\label{eq:l05-normal-entropy}
\end{equation}
\end{proposition}
\begin{proof}
Taking the negative logarithm of~\eqref{eq:l05-normal} times the reference scale and averaging gives
\begin{align}
h_{x_{\mathrm{ref}}}(X)
&=\ln\left(\frac{\sqrt{2\pi}\sigma}{x_{\mathrm{ref}}}\right)
 +\frac{\mathbb E[(X-\mu)^2]}{2\sigma^2} \notag\\
&=\ln\left(\frac{\sqrt{2\pi}\sigma}{x_{\mathrm{ref}}}\right)+\frac12,
\label{eq:l05-normal-entropy-proof}
\end{align}
which yields the stated result because $\frac12 = \frac12\ln\ee$.
\end{proof}
\textit{Significance:} exponentiating entropy yields an effective width proportional to $\sigma$ ($\sqrt{2\pi\ee}\,\sigma \approx 4.13\,\sigma$). Translation by $\mu$ changes neither width nor entropy. Doubling $\sigma$ increases entropy by $\ln2$.

Gaussian parameters can be estimated directly from unbinned data using maximum likelihood:
\begin{align}
\ln\mathcal L(\mu,\sigma^2)
&=-\frac N2\ln(2\pi\sigma^2)
 -\frac1{2\sigma^2}\sum_{n=1}^N(x_n-\mu)^2,\label{eq:l05-normal-likelihood}\\
\frac{\partial\ln\mathcal L}{\partial\mu}
&=\frac1{sigma^2}\sum_n(x_n-\mu)=0
\quad\Longrightarrow\quad\widehat\mu=\overline x,\label{eq:l05-normal-mean-mle}\\
\frac{\partial\ln\mathcal L}{\partial\sigma^2}
&=-\frac N{2\sigma^2}+\frac{\sum_n(x_n-\overline x)^2}{2\sigma^4}=0
\quad\Longrightarrow\quad
\widehat\sigma^2=\frac1N\sum_n(x_n-\overline x)^2.
\label{eq:l05-normal-var-mle}
\end{align}
The estimator with denominator $N$ is the maximum likelihood estimator. Denominator $N-1$ provides the unbiased sample variance. A finite positive width requires non-zero sample spread.

\subsection{A phase lives on a circle}
\label{subsec:l05-vonmises}
Circadian rhythms and cell-cycle phases are periodic: phase 0 and phase $2\pi$ represent identical biological states. A Gaussian on the real line assigns different values to angles that differ by $2\pi$, so it is not directly a circular density. Wrapping a real Gaussian across the circle produces an infinite series. The von Mises distribution provides a closed-form, periodic alternative.

\begin{definition}{Von Mises distribution}{l05-vonmises}
For angular coordinate $\theta\in[0,2\pi)$, mean direction $\mu$ modulo $2\pi$, and dimensionless concentration $\kappa\geq0$,
\begin{equation}
p(\theta)=\frac{\exp[\kappa\cos(\theta-\mu)]}{2\pi I_0(\kappa)},
\qquad
I_0(\kappa)=\frac1{2\pi}\int_0^{2\pi}\exp(\kappa\cos u)\,\dd u.
\label{eq:l05-vonmises}
\end{equation}
Here $I_0$ is the modified Bessel function of the first kind, order zero. The integral definition ensures normalization over the circle.
\end{definition}
Angle $\theta$ is in radians. Increasing $\kappa$ concentrates probability around $\mu$. At $\kappa=0$, the distribution is uniform and $\mu$ is unidentifiable.

For $\kappa\to0$, expanding the exponential shows that $I_0(\kappa)\to1$, yielding
\begin{equation}
p(\theta)\longrightarrow\frac1{2\pi}.
\label{eq:l05-vonmises-uniform}
\end{equation}
For $\kappa\gg1$, let $\delta=\theta-\mu$ near zero and expand the cosine:
\begin{align}
\kappa\cos\delta
&=\kappa-\frac{\kappa\delta^2}{2}+O(\kappa\delta^4),\label{eq:l05-cos-expansion}\\
2\pi I_0(\kappa)
&\simeq\ee^\kappa\int_{-\infty}^{\infty}
\exp(-\kappa\delta^2/2)\,\dd\delta
=\ee^\kappa\sqrt{\frac{2\pi}{\kappa}},\label{eq:l05-bessel-asymptotic}\\
p(\theta)&\simeq\sqrt{\frac{\kappa}{2\pi}}
\exp\left(-\frac{\kappa\delta^2}{2}\right),
\qquad \sigma_{\mathrm{local}}^2\simeq\frac1\kappa.
\label{eq:l05-vonmises-normal}
\end{align}
Typical deviations are $\delta=O(\kappa^{-1/2})$, making the neglected exponent terms $O(1/\kappa)$. This distribution is a local Gaussian for large finite $\kappa$. The exact infinite-concentration limit is a Dirac delta on the circle rather than a Gaussian of finite width.

\begin{example}{Phase clustering across the periodic boundary}{l05-phase-example}
Two phases $\theta_1=0.1$ and $\theta_2=2\pi-0.1$ radians lie adjacent on the circle, though their linear arithmetic mean is $\pi$. Averaging their 2D unit vectors yields zero sine component and positive cosine component, giving circular mean direction zero. A von Mises density with $\kappa=100$ has local standard deviation approximately $0.1$ radians (corresponding to $24 \times 0.1 / (2\pi) \simeq 0.38$ hours for circadian rhythms). A cluster spanning the $2\pi$ wrap represents a single population.
\end{example}

Periodic biological processes, including circadian rhythms and cell division cycles, provide direct physical applications for circular distributions. Different metabolic processes can exhibit distinct phase offsets while sharing a common 24-hour period. Environmental cues, such as light-dark cycles, synchronize coupled biological oscillators. A single von Mises component models unimodal phase clustering, while multimodal or non-sinusoidal rhythms require mixtures or generalized phase models.

\section{Flow cytometry: a biological population becomes a point cloud}
\label{sec:l05-cytometry}
\subsection{What is measured on each cell?}
\label{subsec:l05-instrument}
Immune cell populations exhibit continuous variation in surface protein expression, internal structure, and size. Flow cytometry measures these properties across individual cells, converting a liquid cell suspension into a multivariate single-cell point cloud.

Hydrodynamic focusing aligns cells into a single-file stream through a laser interrogation zone. Aggregates and doublets can be misidentified as large single cells, requiring pulse-geometry gating (comparing pulse height to pulse area). Optical detectors record light scatter and fluorescence. Forward scatter provides an empirical proxy for cell cross-sectional area, while side scatter reflects internal complexity and granularity. Neither measurement constitutes an exact physical power law of cell radius. Scattered intensity depends on cell morphology, refractive index, and optical collection geometry.

Fluorescently conjugated antibodies bind designated target epitopes. Molecular binding is subject to cross-reactivity and non-specific binding, requiring isotype and negative controls. Fluorophores absorb laser light and emit photons across a broad wavelength spectrum. Dichroic mirrors and bandpass filters direct emitted photons to photomultiplier tubes or avalanche photodiodes, where electronic amplification converts photon flux into digital channel values. When detector gain and optical channels operate in linear regimes, recorded intensity is proportional to labeled epitope abundance. Across distinct fluorophore channels, identical numerical readings do not imply identical protein numbers.

Modern instruments measure tens of parameters at acquisition rates exceeding $10^4$ cells per second, generating datasets with $10^5$--$10^6$ single-cell records. Fluorescence-activated cell sorting (FACS) adds electrostatic droplet deflection to sort target cells into collection vessels based on optical triggers. Cells undergo optical interrogation before droplet formation and charging.

\subsection{Gates simplify distributions, sometimes too much}
\label{subsec:l05-gates}
A \emph{gate} defines a decision boundary within the measurement space. For example, a bivariate plot of side scatter versus CD45 expression resolves major leukocyte populations in peripheral blood. CD45 serves as a pan-leukocyte marker, while side scatter separates low-granularity lymphocytes from intermediate monocytes and highly granular granulocytes.

Drawing a polygonal gate on a two-dimensional projection isolates a cellular subpopulation. However, two-dimensional projections discard the remaining coordinates. Overlapping subpopulations in one plane may separate completely in a third marker. Conversely, sharp boundaries can impose artificial dichotomies onto a continuous biological differentiation spectrum. High-dimensional clustering algorithms automate gating, but their stability across biological replicates requires validation.

\subsection{Spectral overlap produces spillover}
\label{subsec:l05-spectra}
Optical filters transmit wavelength intervals rather than specific fluorophore emissions. A ``525/50 nm'' filter centers at 525 nm with a nominal passband of $525\pm25$ nm. A ``585/40 nm'' filter spans $585\pm20$ nm. If two fluorophore emission spectra overlap a single detector passband, both fluorophores generate signal in that channel.

If $s_j(\lambda)$ represents the emission spectrum per unit abundance of fluorophore $j$, and $R_i(\lambda)$ represents the detector response of channel $i$, the linear mixing coefficient is proportional to the spectral overlap integral:
\begin{equation}
M_{ij}\propto\int R_i(\lambda)s_j(\lambda)\,\dd\lambda.
\label{eq:l05-spectral-integral}
\end{equation}
The matrix elements $M_{ij}$ represent these optical overlap integrals. A small fractional spillover from an intensely expressed marker can overwhelm the primary signal of a dim target. For example, a 5\% spillover from a bright signal of 1000 units contributes 50 units, exceeding a primary signal of 10 units. Panel design selects compatible fluorophore-filter combinations to minimize overlap, but cannot eliminate spectral spillover completely.

\section{Spillover creates covariance, and compensation attempts to remove it}
\label{sec:l05-compensation}
\subsection{A deliberately simple two-channel model}
\label{subsec:l05-mixing}
Let $X_0, Y_0$ denote unmixed biological signals in calibrated units, and let $X, Y$ denote recorded detector signals. A symmetric two-channel model demonstrates the underlying algebra by routing a fraction $\alpha$ of each channel into the opposing detector:
\begin{equation}
\begin{pmatrix}X\\Y\end{pmatrix}
=M\begin{pmatrix}X_0\\Y_0\end{pmatrix},
\qquad M=\begin{pmatrix}1-\alpha&\alpha\\\alpha&1-\alpha\end{pmatrix}.
\label{eq:l05-mixing-matrix}
\end{equation}
Equivalently, $X=(1-\alpha)X_0+\alpha Y_0$ and $Y=\alpha X_0+(1-\alpha)Y_0$. For physical detectors, $0\leq\alpha<1/2$, ensuring that each primary detector captures the majority of its target signal.

\begin{figure}[htbp]
\centering
\begin{tikzpicture}[x=1cm,y=1cm,>=stealth]
\node (x0) at (0,1) {$X_0$}; \node (y0) at (0,-1) {$Y_0$};
\node (x) at (6,1) {$X$}; \node (y) at (6,-1) {$Y$};
\draw[->,thick,primary] (x0)--node[above] {$1-\alpha$} (x);
\draw[->,thick,primary] (y0)--node[below] {$1-\alpha$} (y);
\draw[->,accent] (x0)--node[pos=.25,below] {$\alpha$} (y);
\draw[->,accent] (y0)--node[pos=.25,above] {$\alpha$} (x);
\end{tikzpicture}
\caption{Schematic symmetric spillover model. Each physical detector receives photons from both fluorophores. The matrix $M$ contains the fractional transmission coefficients.}
\label{fig:l05-spillover}
\end{figure}

To evaluate covariance induced purely by measurement, assume $X_0$ and $Y_0$ are uncorrelated biological signals with variances $s_x^2$ and $s_y^2$. Let $u=X_0-\mathbb E[X_0]$ and $v=Y_0-\mathbb E[Y_0]$ denote centered fluctuations, satisfying $\mathbb E[u]=\mathbb E[v]=\mathbb E[uv]=0$. By linearity:
\begin{equation}
\delta X=(1-\alpha)u+\alpha v,\qquad
\delta Y=\alpha u+(1-\alpha)v.
\label{eq:l05-centered-mixing}
\end{equation}
Squaring and evaluating expectations yields:
\begin{align}
\operatorname{Var}(X)
&=(1-\alpha)^2\mathbb E[u^2]
+2\alpha(1-\alpha)\mathbb E[uv]+\alpha^2\mathbb E[v^2] \notag\\
&=(1-\alpha)^2s_x^2+\alpha^2s_y^2,\label{eq:l05-mixed-varx}\\
\operatorname{Var}(Y)&=\alpha^2s_x^2+(1-\alpha)^2s_y^2.
\label{eq:l05-mixed-vary}
\end{align}
Multiplying the centered fluctuations gives the covariance:
\begin{align}
\operatorname{Cov}(X,Y)
&=\mathbb E\big[\alpha(1-\alpha)u^2
+((1-\alpha)^2+\alpha^2)uv+\alpha(1-\alpha)v^2\big] \notag\\
&=\boxed{\alpha(1-\alpha)(s_x^2+s_y^2)}.
\label{eq:l05-mixed-cov}
\end{align}
\textit{Significance:} instrument mixing generates positive covariance between measured channels even when biological expression is completely uncorrelated. At $\alpha=0$, covariance vanishes. Zero initial covariance $\operatorname{Cov}(X_0, Y_0)=0$ suffices to derive~\eqref{eq:l05-mixed-cov}.

\begin{example}{Spurious correlation from independent signals}{l05-spillover-example}
Let $s_x^2=s_y^2=1$ and $\alpha=0.1$. Then:
\begin{equation}
\operatorname{Var}(X)=\operatorname{Var}(Y)=0.9^2+0.1^2=0.82,
\quad\operatorname{Cov}(X,Y)=0.18,
\quad\rho=\frac{0.18}{0.82}=\frac9{41}\simeq0.220.
\label{eq:l05-spillover-numbers}
\end{equation}
An instrument spillover of 10\% generates a false Pearson correlation coefficient of $\rho \approx 0.22$.
\end{example}

\subsection{Calibrated inversion and noise propagation}
\label{subsec:l05-inverse}
\begin{proposition}{Compensation for the symmetric model}{l05-inverse}
For $\alpha\neq1/2$, the unmixed signals in a noiseless system are recovered by matrix inversion:
\begin{equation}
\begin{pmatrix}X_0\\Y_0\end{pmatrix}
=\frac1{1-2\alpha}
\begin{pmatrix}1-\alpha&-\alpha\\-\alpha&1-\alpha\end{pmatrix}
\begin{pmatrix}X\\Y\end{pmatrix}.
\label{eq:l05-compensation}
\end{equation}
\end{proposition}
\begin{proof}
The determinant is $\det M = (1-\alpha)^2-\alpha^2=1-2\alpha$. Computing the classical $2\times2$ matrix inverse yields~\eqref{eq:l05-compensation}. Multiplying $M^{-1}M$ returns the identity matrix.
\end{proof}
\textit{Significance:} compensation uses physical calibration to invert optical mixing. It is not an empirical procedure to force sample covariance to zero.

Summing and differencing the measured coordinates reveals the sensitivity:
\begin{equation}
X+Y=X_0+Y_0,\qquad X-Y=(1-2\alpha)(X_0-Y_0).
\label{eq:l05-sum-difference}
\end{equation}
At $\alpha=1/2$, both detectors record the identical average, and signal difference is lost. Near $\alpha=1/2$, matrix inversion amplifies difference-channel noise by $(1-2\alpha)^{-1}$. If measurement introduces detector noise $\bm{\eta}$, compensated coordinates carry noise $M^{-1}\bm{\eta}$.

Single-stain controls calibrate the matrix columns. Cells stained with fluorophore $X_0$ alone yield $Y/X=\alpha/(1-\alpha)$ after background subtraction, uniquely fixing $\alpha$. In multi-color panels, $d$ single-stain controls determine the full $d\times d$ spillover matrix.

\begin{remark}
Multimodal marker distributions in single-stain controls do not invalidate linear compensation. Matrix inversion applies on an event-by-event basis regardless of population distribution shapes. Artifacts arise from miscalibrated coefficients, uncorrected autofluorescence, detector non-linearity, or fluorophore saturation. Overcompensation induces spurious negative covariance, while undercompensation leaves residual positive covariance. Neither artifact indicates true biological interaction.
\end{remark}

\subsection{Negative compensated values and arcsinh scaling}
\label{subsec:l05-arcsinh}
The matrix inverse~\eqref{eq:l05-compensation} subtracts channel signals. For a cell expressing zero target molecules, Poisson photon counting noise and background subtraction frequently produce negative compensated estimates $\widehat{X}_0 < 0$. Negative compensated estimates represent measurement noise and background subtraction around zero rather than unphysical negative molecule counts.

Standard logarithmic scales cannot display non-positive numbers ($\ln x$ diverges as $x \to 0^+$). Truncating negative values to zero creates an artificial spike at zero and inflates population means. Shifting the dataset alters the biological origin. The standard transformation in single-cell cytometry is the inverse hyperbolic sine with scale factor $c>0$:
\begin{equation}
t(x)=\operatorname{arsinh}(x/c)
=\ln\left[\frac xc+\sqrt{1+\left(\frac xc\right)^2}\right].
\label{eq:l05-arcsinh}
\end{equation}
This function is smooth, odd, monotonic, and defined across all real values. Its derivative and limits are:
\begin{align}
\frac{\dd t}{\dd x}&=\frac1{c\sqrt{1+(x/c)^2}}>0,\label{eq:l05-arcsinh-derivative}\\
t(x)&=\frac xc+O\big((x/c)^3\big),\qquad \abs{x}\ll c,\label{eq:l05-arcsinh-small}\\
t(x)&\simeq\operatorname{sgn}(x)\ln\left(\frac{2\abs{x}}c\right),
\qquad \abs{x}\gg c.
\label{eq:l05-arcsinh-large}
\end{align}
For small signals ($\abs{x}\ll c$), the transform is linear ($t(x) \approx x/c$), preserving symmetric background noise around zero. For large signals ($\abs{x}\gg c$), it compresses dynamic range logarithmically. The scale factor $c$ sets the linear-to-logarithmic crossover. This transformation alters geometric distances and density shapes, which must be considered when interpreting clusters.

\subsection{Mutual information as a calibration diagnostic}
\label{subsec:l05-mi-calibration}
In single-stain control samples, minimizing mutual information between compensated channels can guide matrix calibration. Unlike linear covariance, mutual information detects nonlinear dependency without requiring Gaussian assumptions. In practice, estimating mutual information from finite two-dimensional empirical distributions requires substantial sample sizes.

Calibration criteria apply to single-stain controls where channels should be statistically independent. In biological samples, true correlated expression must be preserved. Optimization must also avoid degenerate transformations that collapse channels to zero variance.

\subsubsection*{Method: diagnosing apparent biological correlation}
\begin{enumerate}
\item Define the physical fluorophores and optical detector configuration before evaluating scatter plots.
\item Measure single-stain controls and unstained cells to calibrate the spillover matrix $M$ and background offsets.
\item Invert the calibrated matrix $M^{-1}$ on single-cell vectors, retaining signed values to preserve noise symmetry.
\item Apply $\operatorname{arsinh}$ scaling with an explicit cofactor $c$ tuned to background noise standard deviations.
\item Confirm that single-stain controls show zero residual covariance, distinguishing true biological co-expression from optical spillover.
\end{enumerate}

\section{The multivariate Gaussian and its information}
\label{sec:l05-multivariate}
\subsection{Covariance fixes orientation and width}
\label{subsec:l05-mvnormal}
Let $\mathbf X\in\mathbb R^k$ be a random measurement vector, $\boldsymbol\mu=\mathbb E[\mathbf X]$ its mean vector, and
\begin{equation}
\Sigma_{ij}=\mathbb E[(X_i-\mu_i)(X_j-\mu_j)]
\label{eq:l05-cov-definition}
\end{equation}
its covariance matrix. The covariance matrix $\Sigma$ is symmetric and positive semidefinite. For any real projection vector $\mathbf v$, $\mathbf v^{\mathsf T}\Sigma\mathbf v = \operatorname{Var}(\mathbf v^{\mathsf T}\mathbf X) \geq 0$. Diagonal entries represent coordinate variances $\sigma_i^2$, while off-diagonal entries represent covariances.

\begin{definition}{Multivariate Gaussian density}{l05-mvnormal}
For positive-definite $\Sigma$, the normalized multivariate Gaussian probability density function is
\begin{equation}
p(\mathbf x)=\frac1{(2\pi)^{k/2}\sqrt{\det\Sigma}}
\exp\left[-\frac12(\mathbf x-\boldsymbol\mu)^{\mathsf T}
\Sigma^{-1}(\mathbf x-\boldsymbol\mu)\right].
\label{eq:l05-mvnormal}
\end{equation}
\end{definition}
The matrix inverse $\Sigma^{-1}$ generalizes $1/\sigma^2$ to $k$ dimensions. The quadratic Mahalanobis distance measures displacement in units of distribution variance, rendering the exponent dimensionless. If coordinate $i$ carries physical unit $u_i$, the joint density carries units $\prod_i u_i^{-1}$. A singular covariance matrix defines a Gaussian supported on a lower-dimensional subspace. The full $k$-dimensional density~\eqref{eq:l05-mvnormal} does not exist on that subspace.

For two variables with positive standard deviations $\sigma_x,\sigma_y$, define the dimensionless Pearson correlation coefficient $\rho=\operatorname{Cov}(X,Y)/(\sigma_x\sigma_y)$. Then:
\begin{align}
\Sigma&=\begin{pmatrix}\sigma_x^2&\rho\sigma_x\sigma_y\\
\rho\sigma_x\sigma_y&\sigma_y^2\end{pmatrix},
\quad \det\Sigma=\sigma_x^2\sigma_y^2(1-\rho^2),\label{eq:l05-bivariate-cov}\\
\Sigma^{-1}&=\frac1{1-\rho^2}
\begin{pmatrix}\sigma_x^{-2}&-\rho/(\sigma_x\sigma_y)\\
-\rho/(\sigma_x\sigma_y)&\sigma_y^{-2}\end{pmatrix}.
\label{eq:l05-bivariate-inverse}
\end{align}
Positive definiteness requires $\abs\rho<1$. With standardized variables $u=(x-\mu_x)/\sigma_x$ and $v=(y-\mu_y)/\sigma_y$:
\begin{equation}
p(x,y)=\frac{\exp\left[-\dfrac{u^2-2\rho uv+v^2}{2(1-\rho^2)}\right]}
{2\pi\sigma_x\sigma_y\sqrt{1-\rho^2}}.
\label{eq:l05-bivariate-density}
\end{equation}
Non-zero correlation $\rho\neq0$ tilts the elliptical contours of equal density. At $\rho=0$, the joint density factorizes into the product of marginal Gaussians $p(x,y)=p(x)p(y)$. For non-Gaussian distributions, zero covariance does not guarantee independence.

\subsection{Gaussian entropy and mutual information}
\label{subsec:l05-gaussian-mi}
Evaluating the expectation of the negative log-density for~\eqref{eq:l05-mvnormal} with $\mathbf q=\mathbf X-\boldsymbol\mu$:
\begin{align}
\mathbb E[\mathbf q^{\mathsf T}\Sigma^{-1}\mathbf q]
&=\operatorname{tr}\left(\Sigma^{-1}\mathbb E[\mathbf q\mathbf q^{\mathsf T}]\right)
=\operatorname{tr}(\Sigma^{-1}\Sigma)=\operatorname{tr}(I_k)=k,\label{eq:l05-quadratic-expectation}\\
h(\mathbf X)&=\frac k2\ln(2\pi)+\frac12\ln\det\Sigma+\frac k2
=\frac12\ln\left[(2\pi\ee)^k\det\Sigma\right].
\label{eq:l05-mv-entropy}
\end{align}
For dimensional coordinates, $\det\Sigma$ is normalized by reference scales $\prod_i x_{\mathrm{ref},i}^2$.

\begin{proposition}{Mutual information of a bivariate Gaussian}{l05-gaussian-mi}
For a jointly Gaussian pair with correlation $\rho$:
\begin{equation}
\boxed{I(X;Y)=-\frac12\ln(1-\rho^2)}.
\label{eq:l05-gaussian-mi}
\end{equation}
\end{proposition}
\begin{proof}
Using $I(X;Y)=h(X)+h(Y)-h(X,Y)$ and determinant $\det\Sigma=\sigma_x^2\sigma_y^2(1-\rho^2)$:
\begin{align}
I(X;Y)
&=\frac12\ln(2\pi\ee\sigma_x^2)
 +\frac12\ln(2\pi\ee\sigma_y^2)
 -\frac12\ln\left[(2\pi\ee)^2\sigma_x^2\sigma_y^2(1-\rho^2)\right] \notag\\
&=-\frac12\ln(1-\rho^2).
\label{eq:l05-mi-cancellation}
\end{align}
All scale constants and marginal variances cancel exactly.
\end{proof}
\textit{Significance:} for jointly Gaussian variables, correlation completely specifies mutual information. The sign of $\rho$ does not alter $I(X;Y)$. At $\rho=0$, mutual information vanishes. As $\abs\rho\to1$, information diverges because a noiseless observation of $X$ determines $Y$ with infinite precision. Finite instrument resolution and detector noise bound this divergence in physical measurements.

\begin{example}{Spurious information induced by optical spillover}{l05-mi-example}
If independent biological inputs $X_0, Y_0$ in Example~\ref{ex:l05-spillover-example} are Gaussian, recorded signals are jointly Gaussian. For $\rho=9/41$:
\begin{equation}
I(X;Y)=-\frac12\ln\left(1-\frac{81}{1681}\right)
=-\frac12\ln\left(\frac{1600}{1681}\right)\simeq0.0247\ \text{nats}.
\label{eq:l05-mi-numeric}
\end{equation}
This mutual information arises purely from optical spillover. For non-Gaussian biological signals, the linear covariance remains identical, but the mutual information depends on higher-order moments.
\end{example}

\section{Principal axes, whitening, and dimension reduction}
\label{sec:l05-pca}
\subsection{Diagonalizing the covariance matrix}
\label{subsec:l05-diagonalization}
A real symmetric covariance matrix $\Sigma$ possesses an orthonormal basis of eigenvectors $\mathbf p_i$ with real eigenvalues $\lambda_i\geq0$. Assembling eigenvectors into the orthogonal matrix $P=(\mathbf p_1 \,|\, \dots \,|\, \mathbf p_k)$:
\begin{equation}
\Sigma=P\Lambda P^{\mathsf T},\qquad
P^{\mathsf T}P=I_k,\qquad
\Lambda=\operatorname{diag}(\lambda_1,\ldots,\lambda_k).
\label{eq:l05-eigendecomposition}
\end{equation}
Centering and projecting data onto these principal axes defines principal coordinates:
\begin{equation}
\mathbf Y=P^{\mathsf T}(\mathbf X-\boldsymbol\mu),\qquad
Y_i=\mathbf p_i^{\mathsf T}(\mathbf X-\boldsymbol\mu).
\label{eq:l05-pc-coordinates}
\end{equation}
Because linear transformations transform covariance as $\operatorname{Cov}(A\mathbf X)=A\Sigma A^{\mathsf T}$:
\begin{equation}
\operatorname{Cov}(\mathbf Y)
=P^{\mathsf T}\Sigma P
=P^{\mathsf T}(P\Lambda P^{\mathsf T})P=\Lambda.
\label{eq:l05-pc-covariance}
\end{equation}
Principal coordinates have zero mutual covariance ($\operatorname{Cov}(Y_i, Y_j)=0$ for $i\neq j$). The variance along axis $i$ is $\lambda_i$. The standard deviation along principal axis $i$ is $\sqrt{\lambda_i}$. An orthogonal basis transformation may include reflections depending on eigenvector sign conventions, which leaves physical and statistical conclusions unchanged.

For positive eigenvalues, $\Sigma^{-1}=P\Lambda^{-1}P^{\mathsf T}$. Substituting $\mathbf x-\boldsymbol\mu=P\mathbf y$ into the Gaussian exponent yields:
\begin{equation}
(\mathbf x-\boldsymbol\mu)^{\mathsf T}\Sigma^{-1}(\mathbf x-\boldsymbol\mu)
=\mathbf y^{\mathsf T}\Lambda^{-1}\mathbf y
=\sum_i\frac{y_i^2}{\lambda_i}.
\label{eq:l05-pca-quadratic}
\end{equation}
Because $\det P = 1$ and $\det\Sigma = \prod_i \lambda_i$, the joint density factorizes into independent Gaussians:
\begin{equation}
p_{\mathbf Y}(\mathbf y)
=\prod_{i=1}^k\frac1{\sqrt{2\pi\lambda_i}}
\exp\left(-\frac{y_i^2}{2\lambda_i}\right).
\label{eq:l05-pca-factorization}
\end{equation}

\begin{remark}
PCA applies to any empirical covariance matrix without assuming Gaussianity. It guarantees uncorrelated principal coordinates. Statistical independence requires additional distributional structure, such as joint Gaussianity.
\end{remark}

\begin{figure}[htbp]
\centering
\begin{tikzpicture}[x=.85cm,y=.85cm,>=stealth]
\draw[->] (-2.8,0)--(3.1,0) node[right] {$x-\mu_x$};
\draw[->] (0,-2.5)--(0,2.8) node[above] {$y-\mu_y$};
\draw[primary,thick,rotate=45] (0,0) ellipse (2.25 and .85);
\draw[primary!50,rotate=45] (0,0) ellipse (1.5 and .57);
\draw[->,very thick,accent] (-1.9,-1.9)--(2.1,2.1) node[above right] {$\mathbf p_1$};
\draw[->,thick,accent] (1.1,-1.1)--(-1.25,1.25) node[above left] {$\mathbf p_2$};
\fill (0,0) circle (2pt);
\end{tikzpicture}
\caption{Schematic equal-density contours of a centered bivariate Gaussian with positive covariance. Principal axes $\mathbf p_1, \mathbf p_2$ align with major and minor ellipse axes. PCA rotates to these axes; whitening scales them to unit circle contours.}
\label{fig:l05-pca-axes}
\end{figure}

\subsection{Whitening: standardizing coordinate variance}
\label{subsec:l05-whitening}
\begin{definition}{PCA whitening}{l05-whitening}
For strictly positive eigenvalues $\lambda_i>0$, the whitened coordinates are defined by
\begin{equation}
\mathbf Z=\Lambda^{-1/2}P^{\mathsf T}(\mathbf X-\boldsymbol\mu).
\label{eq:l05-whitening}
\end{equation}
The covariance matrix of $\mathbf Z$ is the identity matrix:
\begin{equation}
\operatorname{Cov}(\mathbf Z)
=\Lambda^{-1/2}\Lambda\Lambda^{-1/2}=I_k.
\label{eq:l05-whitened-cov}
\end{equation}
\end{definition}
\textit{Significance:} PCA rotates coordinates to align with ellipse axes. Whitening rescales each principal coordinate by $\lambda_i^{-1/2}$, producing unit variances and spherical Gaussian contours. Whitened coordinates are dimensionless. For singular matrices, whitening applies only across positive-eigenvalue subspaces. Small eigenvalues amplify noise and destabilize whitening.

These operations address distinct goals. Physical compensation inverts optical detector spillover using calibration controls. PCA identifies orthogonal axes of observed variance. Whitening standardizes variance across components. Observed sample covariance alone cannot distinguish biological co-regulation from optical spillover. Whitening a sample with genuine biological correlation does not recover unmixed biological signals.

\subsection{Dimensionality reduction and explained variance}
\label{subsec:l05-truncation}
Order eigenvalues by descending magnitude: $\lambda_1\geq\lambda_2\geq\cdots\geq\lambda_k\geq0$. The total sample variance equals the matrix trace:
\begin{equation}
\operatorname{tr}\Sigma=\sum_{i=1}^k\lambda_i,
\qquad
f_i=\frac{\lambda_i}{\sum_{j=1}^k\lambda_j},\qquad
f_{1:m}=\frac{\sum_{i=1}^m\lambda_i}{\sum_{j=1}^k\lambda_j}.
\label{eq:l05-explained-variance}
\end{equation}
Here $f_i$ represents the fractional variance explained by component $i$, and $f_{1:m}$ represents the cumulative fraction captured by the first $m$ components. Retaining $m < k$ components yields the low-rank reconstruction:
\begin{equation}
\widehat{\mathbf x}=\boldsymbol\mu+\sum_{i=1}^m y_i\mathbf p_i.
\label{eq:l05-reconstruction}
\end{equation}
By eigenvector orthogonality, the expected squared reconstruction error equals the discarded eigenvalues:
\begin{equation}
\mathbb E\!\left[\norm{\mathbf X-\widehat{\mathbf X}}^2\right]
=\mathbb E\!\left[\sum_{i=m+1}^kY_i^2\right]
=\sum_{i=m+1}^k\lambda_i.
\label{eq:l05-discarded-variance}
\end{equation}
An orthogonal rotation preserves total variance. Linear projection discards the variance in omitted directions ($\sum_{i=m+1}^k \lambda_i$). PCA provides an optimal low-rank approximation under squared error.

\begin{example}{Two-marker principal component analysis}{l05-pca-example}
Consider centered, scaled fluorescent marker signals with covariance:
\begin{equation}
\Sigma=\begin{pmatrix}4&3\\3&4\end{pmatrix}.
\label{eq:l05-pca-example-matrix}
\end{equation}
Solving $\det(\Sigma-\lambda I)=(4-\lambda)^2-9=0$ yields eigenvalues $\lambda_1=7$ and $\lambda_2=1$. Normalized eigenvectors and principal coordinates are:
\begin{equation}
\mathbf p_1=\frac1{\sqrt2}\begin{pmatrix}1\\1\end{pmatrix},\qquad
\mathbf p_2=\frac1{\sqrt2}\begin{pmatrix}1\\-1\end{pmatrix},\qquad
Y_1=\frac{\delta X+\delta Y}{\sqrt2},\quad
Y_2=\frac{\delta X-\delta Y}{\sqrt2}.
\label{eq:l05-pca-example-vectors}
\end{equation}
Component $Y_1$ captures $7/(7+1)=87.5\%$ of the total variance. Reconstructing data with $Y_1$ alone sets $Y_2=0$, giving an expected squared reconstruction error equal to $\lambda_2=1$. Whitening retains both directions, scaling them as $Z_1=Y_1/\sqrt7$ and $Z_2=Y_2/\sqrt1$ to achieve unit variance.
\end{example}

\subsubsection*{Method: executing PCA on quantitative biological data}
\begin{enumerate}
\item Define biological cohorts and quality-filter measurement channels; apply compensation to eliminate instrument spillover.
\item Explicitly standardize units across features; mean-center all features. Scaling choices alter PCA eigenvectors.
\item Calculate sample covariance $\Sigma$, compute its eigendecomposition, and rank components by descending eigenvalues $\lambda_i$.
\item Project single-cell coordinates onto leading eigenvectors and record explained-variance fractions $f_{1:m}$.
\item Evaluate whether discarded dimensions contain biological variation rather than measurement noise.
\end{enumerate}
In high-dimensional transcriptomics ($k \sim 10^4$), leading components capturing 30\% of variance indicate concentrated coordinated expression programs. Conversely, in low dimensions ($k=3$), a component capturing 33\% reflects isotropic variation. High variance does not guarantee biological importance: technical batch effects and cell-cycle state often dominate leading components, while critical cell subtypes separate along minor axes.

\section{Nonlinear structure and distance-based ordination}
\label{sec:l05-limitations}
A curved or horseshoe-shaped point cloud may possess an intrinsic one-dimensional manifold structure, yet linear rotation cannot unbend the curve. PCA coordinates remain uncorrelated while exhibiting deterministic nonlinear dependence. For example, if $X$ is symmetric about zero and $Y=X^2$, then $\operatorname{Cov}(X,Y)=\mathbb E[X^3]-\mathbb E[X]\mathbb E[X^2]=0$, even though $Y$ is deterministically determined by $X$. Covariance tests linear correlation rather than general statistical independence.

Ecological and single-cell comparisons often provide natural pairwise distance matrices, such as UniFrac or Jaccard distances, without supplying Euclidean feature vectors. Principal coordinates analysis (PCoA, or classical multidimensional scaling) operates directly on these pairwise dissimilarities. It differs from applying PCA to raw feature coordinates.

Given an $n\times n$ symmetric dissimilarity matrix $D$ with zero diagonal ($D_{ii}=0$), let $D^{(2)}$ denote the element-wise squared distance matrix ($D^{(2)}_{ij}=D_{ij}^2$). Define the centering matrix $J=I_n-\frac1n\mathbf1\mathbf1^{\mathsf T}$, where $\mathbf1=(1,\dots,1)^{\mathsf T}$. The centered Gram matrix is
\begin{equation}
B=-\frac12 JD^{(2)}J.
\label{eq:l05-pcoa}
\end{equation}
For Euclidean distances, $B$ equals the matrix of centered inner products: $B_{ij}=(\mathbf r_i-\overline{\mathbf r})\cdot(\mathbf r_j-\overline{\mathbf r})$. Eigendecomposition $B=V\Theta V^{\mathsf T}$ yields embedded Euclidean coordinates $\mathbf z_i = \Theta^{1/2}\mathbf v_i$ for positive eigenvalues $\theta_i>0$.

When dissimilarities violate Euclidean metric properties, the Gram matrix $B$ develops negative eigenvalues. Negative eigenvalues indicate that the observed pairwise distances cannot be embedded in Euclidean space without geometric distortion. Nonlinear manifold learning methods (such as ISOMAP, t-SNE, and UMAP) address curved geometries, while eigenworm analysis applies PCA to postural dynamics in subsequent modules.

\section*{Summary and exam checks}
\begin{itemize}
\item \textbf{ANI and UniFrac:} average nucleotide identity requires alignment coverage to define microbial species bounds. Unweighted UniFrac measures the fraction of evolutionary branch length unique to either community; weighted UniFrac incorporates relative lineage abundances.
\item \textbf{PDF versus CDF:} a probability density function carries units of inverse coordinate ($1/[x]$) and represents probability per interval width; its point value can exceed unity. The empirical cumulative distribution function (ECDF) provides a bin-free representation that converges to the true CDF.
\item \textbf{Differential entropy and mutual information:} differential entropy shifts with coordinate scaling ($h(aX)=h(X)+\ln|a|$) and requires an explicit reference scale. Mutual information $I(X;Y)$ is dimensionless, non-negative, and invariant under separate smooth invertible transformations.
\item \textbf{Gaussian and von Mises distributions:} the Gaussian distribution maximizes differential entropy for a fixed variance, with effective width $\sqrt{2\pi\ee}\,\sigma$. The von Mises distribution models circular phase data, approaching a uniform distribution at $\kappa \to 0$ and a Gaussian with $\sigma^2 \simeq 1/\kappa$ as $\kappa \to \infty$.
\item \textbf{Fluorescence spillover and compensation:} optical spectral overlap introduces spurious positive covariance between independent biological markers. Calibrated matrix inversion $M^{-1}$ unmixes signals but amplifies detector noise. Negative compensated values represent symmetric background noise and must be retained using $\operatorname{arsinh}$ scaling.
\item \textbf{PCA and whitening:} diagonalizing covariance rotates coordinates to orthogonal principal axes. PCA rotation preserves total variance, while whitening scales axes to unit variance. Dimensionality reduction discards minor axes, minimizing mean squared reconstruction error. Uncorrelated principal coordinates imply statistical independence only under joint Gaussianity.
\end{itemize}

\section*{Further Reading}
\begin{itemize}
\item \textit{Physical Models of Biological Systems}, \texttt{models-notes-2026.pdf}, Sections 4.6--4.7: ANI and UniFrac; Sections 5.1--5.3: densities, differential entropy, Gaussian and von Mises distributions; Sections 5.4--5.6: flow cytometry, compensation, multivariate Gaussians, and principal components.
\item Philip Nelson, \textit{Physical Models of Living Systems}, course-supplied PENN edition: probabilistic modeling and quantitative single-cell biological measurements.
\end{itemize}
